% !PRINTED 529
\weberpage{529}
%<<CONT
als eine zu $Q_1$ gehörige Function, so sind die beiden dazu conjugirten Werthe
\weq{7}{\begin{aligned}
y_1 &= (\alpha-\alpha_2)(\alpha_3-\alpha_1)\\
y_2 &= (\alpha-\alpha_3)(\alpha_1-\alpha_2).
\end{aligned}}
Diese drei Grössen sind die Wurzeln einer cubischen Partialresolvente, deren Coëfficienten als Invarianten von $f$ leicht zu bilden sind (vergl. \S.\,64, 65, wo die Grössen $y$, $y_1$, $y_2$ mit $U$, $V$, $W$ bezeichnet waren). Die Gleichung wird
\weq{8}{y^3 - Ay + \sqrt{D} = 0.}
Ihre Discriminante ist $B^2$, und für das Product der Wurzeldifferenzen kann man das Vorzeichen leicht (durch Vergleichung der Glieder höchster Ordnung nach der \gesperrt{Gauss}'schen Definition, \S.\,45, oder auch aus \S.\,64) bestimmen. So erhält man:
\weq{9}{(y-y_1)(y-y_2)(y_1-y_2) = B.}
Die Gleichung (8) ist, wie schon bemerkt, eine Normalgleichung, und es ist leicht, $y_1$, $y_2$ durch $y$ wirklich darzustellen. Wenn man (9) mit $(y_1-y_2)$ multiplicirt, so findet man durch Benutzung der Formeln $(y-y_1)(y-y_2) = 3y^2 - A$, $(y_1-y_2)^2 = y^2 - 4y_1y_2 = 4A - 3y^2$:
\weq{10}{
\begin{aligned}
B(y_1-y_2) &= -9y^4 + 15Ay^2 - 4A^2\\
 &= 6Ay^2 + 9y\sqrt{D} - 4A^2\\
y_1+y_2 &= -y,
\end{aligned}
}
wodurch $y_1$, $y_2$ bestimmt sind. Adjungirt man also fernerhin $y$, so reducirt sich die Gruppe der biquadratischen Gleichung auf
\[
Q_1 = 1, \ (0,2)(1,3), \ (0,1)(2,3), \ (0,3)(1,2).
\]
Diese Gruppe ist imprimitiv (und zwar nach drei Arten), und die Gleichung kann jetzt durch zwei Quadratwurzeln gelöst werden. Am besten gelangt man dazu auf folgendem Wege. Die drei Grössen
\weq{11}{
\begin{aligned}
v_1 &= (\alpha+\alpha_1-\alpha_2-\alpha_3)^2\\
v_2 &= (\alpha-\alpha_1+\alpha_2-\alpha_3)^2\\
v_3 &= (\alpha-\alpha_1-\alpha_2+\alpha_3)^2
\end{aligned}
}
gestatten alle die Permutationen der Gruppe $Q_1$ und sind daher rational durch $y$ darstellbar. Man erhält diese Darstellung folgendermaassen: Es ist
%CONT>>
